% Generated draft: requires mathematical and source review.
\documentclass[11pt]{article}
\usepackage[margin=25mm]{geometry}
\usepackage{fontspec}
\setmainfont{Latin Modern Roman}
\usepackage{amsmath,amssymb,mathtools,unicode-math}
\setmathfont{Latin Modern Math}
\usepackage{enumitem,needspace,xurl,hyperref}
\setlength{\parindent}{0pt}
\setlength{\parskip}{5pt}
\setlength{\emergencystretch}{3em}
\newcommand{\N}{\mathbb{N}}
\newcommand{\Z}{\mathbb{Z}}
\newcommand{\Q}{\mathbb{Q}}
\newcommand{\R}{\mathbb{R}}
\newcommand{\C}{\mathbb{C}}
\newcommand{\F}{\mathbb{F}}
\newcommand{\A}{\mathbb{A}}
\newcommand{\PP}{\mathbb{P}}
\newcommand{\E}{\mathbb{E}}
\newcommand{\eps}{\varepsilon}
\newcommand{\dd}{\,\mathrm{d}}
\begin{document}
\section*{UnsolvedMath ID: 800693}
\textbf{Source problem code:} catalogue entry
\subsection{Definitions and mathematical statement}
Let \(F\) be a field, let \(F^{\mathrm{sep}}\) be a separable closure, and put \(G_F=\operatorname{Gal}(F^{\mathrm{sep}}/F)\). Under the ordinary absolute Galois cohomology convention, \(\operatorname{cd}(F)\leq 2\) means that \(H^q(G_F,M)=0\) for every discrete torsion \(G_F\)-module \(M\) and every integer \(q>2\). The supplied record imposes no restriction on the characteristic of \(F\), and does not specify a different cohomological dimension convention.

A central simple \(F\)-algebra is a finite-dimensional associative unital \(F\)-algebra with center \(F\) and no nonzero proper two-sided ideals. Its degree is \(\deg(A)=\sqrt{\dim_F A}\). Every such algebra has the form \(A\cong M_r(D)\), where \(D\) is a central division \(F\)-algebra, unique up to \(F\)-isomorphism. A central division algebra \(D\) is a crossed product if it contains a subfield \(K\) such that \(K/F\) is finite Galois and \([K:F]=\deg(D)\); such a \(K\) is a maximal subfield.

The supplied question is: if \(\operatorname{cd}(F)\leq 2\), is every central simple \(F\)-algebra a crossed product? Under the convention that this property is tested on the division algebra representative, its precise formulation is
\[
\begin{gathered}
\text{For every field }F\text{ with }\operatorname{cd}(F)\leq 2
\text{ and every finite-dimensional central division }F\text{-algebra }D,\\
\text{there exists a subfield }K\subseteq D\text{ such that }
K/F\text{ is finite Galois and }[K:F]=\deg(D).
\end{gathered}
\]
The qualification concerning the division algebra representative is unresolved in the supplied excerpt. If a crossed product central simple algebra instead means that \(A\) itself contains a finite Galois field extension \(K/F\) of degree \(\deg(A)\), the corresponding question quantifies over \(A\) itself and requires \([K:F]=\deg(A)\). These conventions must not be identified without checking the source's definitions; the displayed formulation is conditional on the division algebra convention.

The source motivates the question by reporting that all examples of non-crossed product algebras known in its 2010 context occur over fields of cohomological dimension at least \(3\). This is supplied historical motivation, not an independently verified assertion about the current literature.
\subsection{Short English statement}
The question asks whether cohomological dimension at most \(2\) forces central simple algebras to be crossed products. Under the division algebra convention, this means that every central division algebra over such a field has a maximal subfield Galois over its center. The intended convention for general central simple algebras, and the cohomological dimension convention in positive characteristic, require source verification.
\subsection{Sources}
\begin{itemize}
\item Local UnsolvedMath registry: permanent ID \(800693\), source problem code \texttt{catalogue entry}, title \emph{Open Problems on Central Simple Algebras --- Problem 3}. The supplied record labels the problem \texttt{open}; this is a registry status, not an independently verified mathematical assessment. Upstream revision: \texttt{d371ee917ac9647fbb06cd6fc529deec6f3dda82}.
\item Auel, Brussel, Garibaldi, and Vishne, \emph{Open Problems on Central Simple Algebras} (2010), Problem 3. Supplied source URL: \url{https://arxiv.org/abs/1006.3304}. The registry attributes the problem to these authors and records proposed year 2010.
\item The supplied provenance reports access on 2026-10-16 and extraction by \texttt{arxiv\_latex}. These are supplied provenance details; the source has not been independently visited or verified for this draft. The statement about all known non-crossed product examples is retained as historical source motivation.
\end{itemize}
\end{document}
